La cuarta revolución industrial se caracteriza por una tendencia hacia la automatización total de la manufactura \cite{Perasso2016}, impulsada por tecnologías como la inteligencia artificial (IA), el internet de las cosas (IoT) y la robótica avanzada. En este contexto, es evidente la premisa sobre la desaparición de muchos y la configuración de nuevos empleos que requieren una alta cualificación y que además son más difíciles de mecanizar. 
\\ \\
Frente al reto latente de esta nueva revolución, para alcanzar el éxito profesional en un mercado laboral cada vez más competitivo y cambiante, es necesario complementar los conocimientos técnicos o científicos con habilidades blandas (habilidades relacionadas con comportamientos, interacción social y la personalidad), las cuales brindan ventaja competitiva, debido a que las empresas están priorizando los perfiles profesionales que cuenten con un mejor repertorio de estas, es decir, que puedan adaptarse a absolutamente todas las necesidades de una organización, puedan trabajar en equipo, puedan ser líderes,  puedan actuar en la resolución de problemas,  puedan ser creativos y proactivos, entre otras cualidades \cite{GonzalezMolina2021}.
\\ \\
Los ingenieros de sistemas no son ajenos a los cambios y desafíos que plantea el panorama laboral actual. Sin embargo, los programas de educación superior en ingeniería de sistemas priorizan la enseñanza de conceptos teóricos o técnicos como la algoritmia, las bases de datos el dominio en lenguajes de programación, entre otros, sobre el desarrollo de competencias interdisciplinarias que permita la integración y aplicación de esos conocimientos. Esto se puede notar en los pensum de ingeniería de sistemas de tanto universidades públicas como privadas, donde el enfoque sigue siendo mayoritariamente técnico y teórico, dejando en segundo plano las habilidades blandas y competencias transversales necesarias para enfrentar los desafíos complejos del mundo laboral contemporáneo \cite{javeriana2021}, \cite{unal2024}.
\\ \\
Ante este panorama, se hace necesario implementar estrategias pedagógicas que promuevan el desarrollo y fortalecimiento de las habilidades blandas en los estudiantes y profesionales de Ingeniería de Sistemas, brindando oportunidades continuas para adquirir, practicar y perfeccionar estas competencias. Una de las estrategias pedagógicas emergentes que ha demostrado gran potencial es la gamificación, la cual aplica elementos y dinámicas propias del juego en contextos educativos, con el propósito de fomentar la motivación, la participación y el aprendizaje significativo \cite{LozadaAvila2017}.
\\ \\
El presente trabajo de grado se enmarca en las iniciativas orientadas al fortalecimiento de las habilidades blandas en los futuros ingenieros de sistemas. Su propósito es diseñar y desarrollar los módulos “Aprender a Aprender” y “Escuch Activa” para el prototipo de plataforma gamificada enfocada en la enseñanza de habilidades blandas, contribuyendo al desarrollo de competencias transversales esenciales para el ejercicio profesional.
\newpage